\section{Algoritmo Branch-and-Bound: Implementação e Auditoria}
\label{sec:bnb}

O algoritmo exato baseado em \textit{Branch-and-Bound} (B\&B) garante a obtenção do ciclo Hamiltoniano de custo mínimo global, explorando implicitamente o espaço fatorial de soluções por meio de podas admissíveis. Nesta seção, avaliamos o comportamento estrutural, o \textit{profiling} por sub-rotina e os gargalos de tempo e memória do B\&B em sua configuração base (\textit{cold start}, inicializado com a rota trivial $1 \to 2 \to \dots \to N \to 1$).

\subsection{Implementação, Estruturas de Dados e Auditoria do Código}

O desenvolvimento do B\&B passou por um processo iterativo de auditoria e refatoração sobre versões geradas inicialmente com auxílio de LLM, cujo código bruto era excessivamente complexo e verboso (mais de $1.300$ linhas, combinando dois motores de busca no mesmo arquivo e gerando sobrecarga estrutural). As principais decisões de engenharia e refatorações aplicadas foram:

\begin{itemize}
    \item \textbf{Estratégia de Busca (DFS vs.\ Best-First):} Adotou-se busca em profundidade (\textit{Depth-First Search} -- DFS) por extensão de caminho com \textit{backtracking in-place}, mantendo apenas os vetores \texttt{path} e \texttt{visited} de tamanho $O(N)$, sem cópias de matrizes ou submatrizes por nó. A escolha evita o estouro de memória inerente à fila de prioridade do \textit{Best-First Search} (que cresce exponencialmente em instâncias com $N \ge 20$). Os vizinhos de cada vértice são pré-ordenados por distância crescente para descobrir boas soluções incumbentes cedo.
    \item \textbf{Limite Inferior (\textit{Lower Bound}) via MST:} Para um caminho parcial $(0 \to \dots \to \text{curr})$ de custo acumulado $g$ e conjunto de vértices não visitados $U$, calcula-se o limite inferior inteiro exato:
    \[
        \mathrm{LB} = g + \mathrm{MST}(U) + \min_{u \in U} d(\text{curr}, u) + \min_{u \in U} d(u, 0),
    \]
    utilizando na raiz o limite de 1-árvore (soma das duas menores arestas distintas incidentes ao vértice $0$ somada a $\mathrm{MST}(V \setminus \{0\})$). Antes de invocar o algoritmo de Prim ($O(|U|^2)$), aplica-se um corte rápido $O(|U|)$ com $g + \min d(\text{curr}, u) + \min d(u, 0) \ge \mathrm{UB}$, além de aborto antecipado dentro do próprio Prim caso a soma parcial já iguale ou supere o limite superior corrente.
    \item \textbf{Otimizações de Baixo Nível e Memória Limitada:} A matriz de distâncias do \texttt{NumPy} é convertida uma única vez para listas nativas do Python para eliminar o \textit{overhead} de empacotamento escalar nos laços críticos. Além disso, implementou-se: (i) algoritmo de Prim denso com \textit{swap-remove} sobre buffers pré-alocados (reduzindo pressão sobre o \textit{garbage collector}); (ii) quebra de simetria de reversão via \textit{bitmask} ordenada por proximidade à origem ($\text{rank}_0[v_1] < \text{rank}_0[v_{N-1}]$ em $O(1)$); e (iii) cache limitado de $\mathrm{MST}(U)$ por \textit{bitmask} com teto fixo de $65.536$ entradas (armazenando apenas custos exatos não abortados). Em testes comparativos antes e depois dessas otimizações, o ganho de desempenho foi de $\approx 2\times$ (ex.: $N=22$ reduziu de $4{,}66\text{ s}$ para $2{,}29\text{ s}$).
    \item \textbf{Isolamento do \textit{Overhead} de Instrumentação:} A medição de memória via \texttt{tracemalloc} foi isolada em execuções dedicadas (\texttt{track\_memory=False} nas coletas de tempo), pois sua ativação introduziu uma desaceleração empírica de $1{,}5\times$ a $2{,}0\times$ no tempo de execução (ex.: \texttt{burma14} passou de $0{,}17\text{ s}$ para $0{,}34\text{ s}$, \texttt{gr17} de $9{,}26\text{ s}$ para $13{,}92\text{ s}$ e \texttt{gr24} de $14{,}95\text{ s}$ para $23{,}73\text{ s}$), o que distorceria a modelagem temporal.
\end{itemize}

\subsection{Resultados Experimentais e Profiling de Sub-rotinas}

O algoritmo foi executado sobre as mesmas 14 instâncias da TSPLIB ($N \in [14, 70]$), com limite estrito de $30\text{ min}$ ($1800\text{ s}$) por execução. A Tabela~\ref{tab:bnb_resumo} apresenta os resultados coletados, registrando como \texttt{NA} o tempo e o ótimo das instâncias interrompidas, mas preservando o melhor custo incumbente, a contagem de nós, o perfil das sub-rotinas e o pico de memória medido em rodada dedicada com \texttt{tracemalloc}.

\begin{table}[htbp]
\centering
\caption{Resultados empíricos e \textit{profiling} do Branch-and-Bound (\textit{cold start}, teto de $1800\text{ s}$).}
\label{tab:bnb_resumo}
\resizebox{\columnwidth}{!}{%
\begin{tabular}{l|c|r|r|c|r|r|r|r|r}
\hline
\textbf{Instância} & \textbf{$N$} & \textbf{$\mathrm{LB}_0$} & \textbf{Gap $\mathrm{LB}_0$} & \textbf{Custo} & \textbf{Nós Exp.} & \textbf{Nós Pod.} & \textbf{Tempo (s)} & \textbf{MST (\%)} & \textbf{Pico (KB)} \\ \hline
\texttt{burma14}   & 14 & 2542 & 23.5\% & 3323         & 6.580       & 40.791      & 0.17  & 82.6\% & 206.7   \\
\texttt{ulysses16} & 16 & 4746 & 30.8\% & 6859         & 2.779.753   & 14.690.413  & 35.84 & 74.9\% & 742.3   \\
\texttt{gr17}      & 17 & 1501 & 28.0\% & 2085         & 295.501     & 2.252.691   & 9.26  & 84.2\% & 800.6   \\
\texttt{gr21}      & 21 & 2252 & 16.8\% & 2707         & 12.892      & 137.820     & 1.27  & 92.9\% & 399.6   \\
\texttt{ulysses22} & 22 & 4866 & 30.6\% & NA (7039)    & 111.869.952 & 643.618.054 & NA    & 78.7\% & 208.6   \\
\texttt{gr24}      & 24 & 1081 & 15.0\% & 1272         & 113.897     & 1.560.923   & 14.95 & 94.4\% & 1681.4  \\
\texttt{fri26}     & 26 & 824  & 12.1\% & 937          & 204.860     & 2.997.809   & 30.94 & 95.2\% & 1478.1  \\
\texttt{bays29}    & 29 & 1622 & 19.7\% & NA (2020)    & 9.248.768   & 150.194.086 & NA    & 95.6\% & 1490.5  \\
\texttt{dantzig42} & 42 & 600  & 14.2\% & NA (699)     & 2.392.064   & 64.596.283  & NA    & 98.4\% & 224.7   \\
\texttt{swiss42}   & 42 & 1107 & 13.0\% & NA (1273)    & 4.030.464   & 83.257.277  & NA    & 97.8\% & 798.8   \\
\texttt{att48}     & 48 & 9029 & 15.0\% & NA (11096)   & 6.627.328   & 127.921.529 & NA    & 96.2\% & 505.8   \\
\texttt{eil51}     & 51 & 385  & 9.6\%  & NA (428)     & 3.870.720   & 83.936.988  & NA    & 97.7\% & 356.3   \\
\texttt{berlin52}  & 52 & 6172 & 18.2\% & NA (8021)    & 6.856.704   & 123.861.967 & NA    & 95.9\% & 303.0   \\
\texttt{st70}      & 70 & 574  & 15.0\% & NA (735)     & 5.849.088   & 110.298.352 & NA    & 96.2\% & 1378.3  \\ \hline
\end{tabular}%
}
\end{table}

\subsection{Modelagem Assintótica e Discussão Crítica}

Para modelar o comportamento empírico do B\&B, ajustaram-se regressões não-lineares e multivariadas via \texttt{SciPy} e \texttt{Scikit-Learn} (Tabela~\ref{tab:regressao_bnb} e Figura~\ref{fig:regressao_bnb}).

\begin{table}[htbp]
\centering
\caption{Modelagem de complexidade empírica do Branch-and-Bound.}
\label{tab:regressao_bnb}
\resizebox{\columnwidth}{!}{%
\begin{tabular}{l|c|c|c|c}
\hline
\textbf{Variável Alvo} & \textbf{Modelo / Preditores} & \textbf{Amostra} & \textbf{$R^2$} & \textbf{RMSE} \\ \hline
Tempo Total ($T$)                & $a \cdot e^{bN}$                                    & 6 resolvidas  & $0.1037$ & $12.986\text{ s}$     \\
$\log_{10}(\text{Nós Exp.})$     & $\beta_1 N + \beta_0$                               & 6 resolvidas  & $0.0027$ & $0.8732$              \\
$\log_{10}(\text{Nós Exp.})$     & $\beta_1 N + \beta_2\,\mathrm{Gap}_{LB_0} + \beta_0$ & 6 resolvidas  & $0.9076$ & $0.2659$              \\
$\log_{10}(\text{Tempo Total})$  & $\beta_1 N + \beta_2\,\mathrm{Gap}_{LB_0} + \beta_0$ & 6 resolvidas  & $0.9515$ & $0.1839$              \\
Tempo MST por Nó Exp.\ ($\mu$s)  & $a \cdot N^2 + c$                                   & 14 instâncias & $0.3088$ & $169.02\,\mu\text{s}$ \\ \hline
\end{tabular}%
}
\end{table}

\begin{figure}[htbp]
    \centering
    \includegraphics[width=\columnwidth]{complexidade_empirica_bnb.png}
    \caption{Análise empírica do Branch-and-Bound: tempo total e ponto de colapso, regressão multivariada com $\mathrm{Gap}_{LB_0}$ e participação da sub-rotina MST no tempo de busca.}
    \label{fig:regressao_bnb}
\end{figure}

As principais conclusões extraídas dos dados empíricos são:

\begin{enumerate}
    \item \textbf{Ponto de Colapso e Insuficiência de $N$ Isolado:} O algoritmo resolveu à otimalidade 6 das 14 instâncias (até $N=26$, \texttt{fri26}, em $30{,}94\text{ s}$), entrando em colapso sistemático a partir de $N \ge 29$ (\texttt{bays29} a \texttt{st70}). Contudo, a regressão puramente em função de $N$ obteve ajuste quase nulo ($R^2 = 0{,}1037$ para $ae^{bN}$ e $R^2 = 0{,}0027$ para $\log_{10}$ dos nós), pois instâncias menores como \texttt{ulysses16} ($N=16$, $35{,}84\text{ s}$, $2{,}78\text{M}$ nós) e \texttt{ulysses22} ($N=22$, \textit{timeout} com $111{,}8\text{M}$ nós) foram drasticamente mais difíceis que \texttt{gr21} ($N=21$, $1{,}27\text{ s}$, $12{,}9\text{k}$ nós) e \texttt{fri26} ($N=26$, $30{,}94\text{ s}$).
    \item \textbf{Dominância do Gap do Limite Inferior na Raiz ($\mathrm{Gap}_{LB_0}$):} A discrepância é explicada diretamente pela folga do limite inferior de 1-árvore na raiz ($\mathrm{Gap}_{LB_0} = (\mathrm{OPT} - \mathrm{LB}_0)/\mathrm{OPT}$). Enquanto \texttt{ulysses16} e \texttt{ulysses22} possuem $\mathrm{Gap}_{LB_0} \approx 30{,}8\%$ (impedindo podas em níveis rasos), \texttt{fri26} apresenta apenas $12{,}1\%$. Ao incluir $\mathrm{Gap}_{LB_0}$ no modelo log-linear ($\log_{10} T = 0{,}3852\,N + 0{,}2092\,\mathrm{Gap}_{LB_0} - 11{,}2250$), o coeficiente de determinação salta para $R^2 = 0{,}9515$ ($\mathrm{RMSE} = 0{,}1839$), indicando crescimento exponencial efetivo de $\approx 2{,}43\times$ por vértice adicional e $\approx 1{,}62\times$ por ponto percentual extra no gap da raiz.
    \item \textbf{Gargalo Computacional na Sub-rotina MST:} O \textit{profiling} revela que o cálculo do \textit{lower bound} via Prim responde em média por $91{,}5\%$ do tempo de busca (variando de $74{,}9\%$ em \texttt{ulysses16} até $98{,}4\%$ em \texttt{dantzig42}, e superando $95\%$ em todas as instâncias com $N \ge 26$). O ajuste quadrático do custo médio por nó ($aN^2 + c$) apresentou $R^2 = 0{,}3088$ porque o custo real de cada chamada depende do tamanho $|U|$ do subconjunto restante na profundidade em que as podas ocorrem e da taxa de acerto do cache de MST: em \texttt{dantzig42} (onde a solução trivial já inicia no ótimo $699$), as podas ocorrem em níveis rasos (subconjuntos $|U|$ grandes, custando $740{,}5\,\mu\text{s}$/nó), mas ainda assim a prova de otimalidade estoura os $30\text{ min}$, demonstrando que a dificuldade central do B\&B reside em \textit{certificar} a otimalidade quando o limite inferior tem folga.
    \item \textbf{Estabilidade Espacial:} Em virtude da estratégia DFS \textit{in-place} com buffers pré-alocados e cache limitado de MST ($65.536$ entradas), o pico de memória medido via \texttt{tracemalloc} manteve-se estritamente contido entre $206{,}7\text{ KB}$ (\texttt{burma14}) e $1{,}64\text{ MB}$ ($1.681{,}4\text{ KB}$ em \texttt{gr24}, execução completa até o ótimo), permanecendo em apenas $1{,}35\text{ MB}$ ($1.378{,}3\text{ KB}$) em \texttt{st70} ($N=70$). Os resultados comprovam que a arquitetura elimina gargalos de RAM e o risco de estouro de memória mesmo quando a árvore de busca ultrapassa dezenas de milhões de nós.
\end{enumerate}
